%
% 13-reell.tex -- Rationale Funktionen mit reellen Koeffizienten
%
% (c) 2026 Prof Dr Andreas Müller
%
\subsection{Rationale Funktionen mit reellen Koeffizienten}
Der Fundamentalsatz der Algebra garantiert eine Faktorisierung 
eines Polynoms mit komplexen Koeffizienten in Linearfaktoren
der Form $z-\alpha$, wobei $\alpha$ eine Nullstelle des Polynoms
ist.
\index{Fundamentalsatz der Algebra}%
Dies ist jedoch nur möglich, wenn komplexe Nullstellen zugelassen
werden.
\index{Nullstelle}%
Es ist jedoch einfach, Polynome anzugeben, die über den reellen Zahlen
nicht in Linearfaktoren zerlegt werden können.
Das Polynom $p(x)=x^2+1$ hat keine reellen Nullstellen.
Über den reellen Zahlen ist $p(x)$ ein irreduzibles Polynom.
Über den komplexen Zahlen ist jedoch $p(x)=(x+i)(x-i)$ eine
Faktorisierung in Linearfaktoren.
\index{Faktorisierung}%

%
% Faktorisierung in lineare und quadratische Faktoren
%
\subsubsection{Faktorisierung in lineare und quadratische Faktoren}
Für ein reelles Polynom 
$p(z) = a_0 + a_1z + \dots + a_{n-1}z^{n-1} + a_nz^n$
kann man also eine Faktorisierung in Linearfaktoren
nicht erwarten.
Betrachtet man das Polynom jedoch als ein Polynom mit komplexen
Koeffizienten, dann garantiert der Fundamentalsatz der Algebra eine
Faktorisierung
\[
p(z)
=
a_n
(z-\alpha_1) 
(z-\alpha_2) 
\cdot
\ldots
\cdot
(z-\alpha_n),
\]
wobei die $\alpha_k$ Nullstellen des Polynoms sind.
Es gilt also $p(\alpha_k)=0$ für alle $k$.

Das Polynom $p(x)$ hat reelle Koeffizienten, die $\bar{a}_k=a_k$
für alle $k$ erfüllen.
Daher kann man die Bedingung $p(\alpha_k)=0$ komplex konjugieren
und bekommt
\begin{align*}
0
=
\overline{p(\alpha_k)}
&=
\bar{a}_0
+
\bar{a}_1 \bar{\alpha}_k
+
\bar{a}_2 \bar{\alpha}_k^2
+
\ldots
+
\bar{a}_{n-1} \bar{\alpha}_k^{n-1}
+
\bar{a}_{n} \bar{\alpha}_k^{n}
\\
&=
a_0
+
a_1 \bar{\alpha}_k
+
a_2 \bar{\alpha}_k^2
+
\ldots
+
a_{n-1} \bar{\alpha}_k^{n-1}
+
a_{n} \bar{\alpha}_k^{n}
\\
&=
p(\bar{\alpha}_k).
\end{align*}
Es folgt, dass $\bar{\alpha}_k$ ebenfalls eine Nullstelle ist.
Falls $\alpha_k=\bar{\alpha}_k$ ist, ist $\alpha_k$ eine reelle
Nullstelle.
Falls $\alpha_k\ne \bar{\alpha}_k$ ist, dann ist $\bar{\alpha}_k$
eine der anderen Nullstellen.
Die Nullstellen eines reellen Polynoms sind als entweder reell,
oder wenn sie nicht reell sind, dann treten sie in Paaren von
konjugiert komplexen Nullstellen auf.

Ein Paar von konjugiert komplexen Nullstellen führt auf den Faktor
\begin{align}
(z-\alpha)(z-\bar{\alpha})
&=
z^2 -(\alpha+\bar{\alpha}) z + \alpha\bar{\alpha}.
\label{buch:ratint:bernoulli:reell:faktoren}
\intertext{Die Summe konjugiert komplexer Zahlen ist
\index{konjugiert}%
$(a+bi)+(a-bi) = 2a $, also das Doppelte des Realteils.
Daher ist der zu berechnende Faktor}
&=
z^2 -2\operatorname{Re}(\alpha)\, z + |\alpha|^2.
\notag
\end{align}
Sowohl $\operatorname{Re}\alpha$ und $|\alpha|^2$ sind reell.
Indem man die Paare konjugiert komplexer Zahlen zu quadratischen
Faktoren wie in \eqref{buch:ratint:bernoulli:reell:faktoren}
zusammenfasst, findet man eine Faktorisierung des Polynoms in
der Form
\begin{equation}
p(z)
=
a_n
(z-\alpha_1)
\cdot\ldots\cdot
(z-\alpha_m)
\cdot
(z^2 + \beta_1 z + \gamma_1)
\cdot \ldots \cdot
(z^2 + \beta_l z + \gamma_l),
\label{buch:ratint:bernoulli:eqn:reellefaktorisierung}
\end{equation}
wobei die $\alpha_1,\dots,\alpha_m$ die reellen Nullstellen sind
und die Faktoren $z^2 +\beta_k z +\gamma_k$ Paare von
konjugiert komplexen, aber nicht reellen Lösungen haben.
Damit die rechte Seite den richtigen Grad hat, muss $n=m+2l$ sein.

%
% Reelle Partialbruchzerlegung
%
\subsubsection{Reelle Partialbruchzerlegung}
In der Zerlegung
\eqref{buch:ratint:bernoulli:eqn:reellefaktorisierung}
können einzelne Faktoren auch mehrfach vorkommen.
Die Konstruktion der Partialbruchzerlegung produziert
daher für jeden Faktor Partialnenner in Potenzen nicht
grösser als in der Faktorisierung.
Dies ist im folgenden Satz zusammengefasst.

\begin{satz}
\label{buch:ratint:bernoulli:satz:reelle-partialbruchzerlegung}
Ist der Nenner $q(z)$ der rationalen Funktion $f(z)=p(z)/q(z)$ 
ein Polynom mit einer Faktorisierung der Form
\begin{equation}
q(z)
=
a_n
(z-\alpha_1)^{s_1}
\cdots
(z-\alpha_m)^{s_m}
(z^2+\beta_1 z + \gamma_1)^{r_1}
\cdots
(z^2+\beta_l z + \gamma_l)^{r_l}
\end{equation}
dann gibt es eine eindeutig bestimmte Partialbruchzerlegung der Form
\begin{equation}
\renewcommand{\arraycolsep}{2pt}
%\renewcommand{\arraystretch}{1.7}
\begin{array}{rcclcccl}
f(z)
&=&&
\displaystyle
\frac{A_1^1    }{ z-\alpha_1       }
&+& \ldots &+&
\displaystyle
\frac{A_1^{r_1}}{(z-\alpha_1)^{s_1}}
\\
&&&\hspace*{11pt}\vdots&&&&\hspace*{17pt}\vdots
\\[2pt]
&&+&
\displaystyle
\frac{A_m^{1}  }{ z-\alpha_m       }
&+& \ldots &+&
\displaystyle
\frac{A_m^{s_m}}{(z-\alpha_m)^{s_m}}
\\[10pt]
&&+&
\displaystyle
\frac{B_1^1    z + C_1^1    }{ z+\beta_1 z + \gamma_1       }
&+& \ldots &+&
\displaystyle
\frac{B_1^{r_1}z + C_1^{r_1}}{(z+\beta_1 z + \gamma_1)^{r_1}}
\\
&&&\hspace*{24pt}\vdots&&&&\hspace*{29pt}\vdots
\\[2pt]
&&+&
\displaystyle
\frac{B_l^1    z + C_1^l    }{ z+\beta_l z + \gamma_l       }
&+& \ldots &+&
\displaystyle
\frac{B_l^{r_l}z + C_l^{r_1}}{(z+\beta_l z + \gamma_l)^{r_l}}.
\end{array}
\label{buch:ratint:bernoulli:eqn:reelle-partialbruchzerlegung}
\end{equation}
Darin sind die Zahlen $A_i^k$, $B_i^k$ und $C_i^k$ reelle Zahlen.
\end{satz}

\begin{proof}
Die Form der Partialbruchzerlegung ist aus der früheren Diskussion,
die zu Satz~\ref{buch:analytisch:rational:satz:partialbruchzerlegung}
geführt hat, klar.
Es muss höchstens noch gezeigt werden, wie diese Zerlegung effizient
ermittelt werden kann.

Offenbar sind $s_i$ Zahlen $A_i^k$ zu bestimmen, insgesamt also
$s_1+\dots+s_m$ Unbekannte für Terme zu den Linearfaktoren.
Zu den quadratischen Faktoren sind je $s_i$ Unbekannte
\index{quadratischer Faktor}%
$B_i^k$ und $C_i^k$ zu bestimmen, also insgesamt
$2r_1+\dots+2r_l$ Unbekannte.
Die Gesamtzahl der Unbekannten ist $s_1+\dots+s_m+2r_1+\dots+2r_l=n$.


Multipliziert man die Gleichung
\eqref{buch:ratint:bernoulli:eqn:reelle-partialbruchzerlegung}
mit dem Nenner $q(z)$, dann bleibt auf der linken Seite das Polynom
$p(z)$ und auf der rechten ein Polynom, dessen Koeffizienten
Linearkombinationen der Unbekannten $A_i^k$, $B_i^k$ und $C_i^k$
sind.
Das Zählerpolynom hat Grad $<n$, es enthält also höchstens $n$
Koeffizienten.
Koeffizientenvergleich liefert dann insgesamt $n$ lineare Gleichungen
für die Unbekannten.
Das entstehende lineare Gleichungssystem hat also $n$ Gleichungen für
\index{lineares Gleichungssystem}%
$n$ Unbekannte.
\end{proof}

Die reelle Partialbruchzerlegung erlaubt die Berechnung einer
Stammfunktion durch Integration der einzelnen Terme.
Für die linearen Terme ist bereits die Stammfunktion
\[
\int\frac{dz}{z-\alpha}
=
\log (z-\alpha)
\]
bekannt.
Die Stammfunktion für die Terme zu den quadratischen Faktoren
des Nenners wird im nächsten Abschnitt bestimmt.

%
% Stammfunktionen für quadratische Faktoren
%
\subsubsection{Stammfunktionen für quadratische Faktoren}
Möchte man das Integral nur im Reellen berechnen, muss man auch 
Partialbrüche der Form
\begin{equation}
f(z)
=
\frac{Az+B}{(z^2 +\beta z + \gamma)^m}
\qquad\text{mit $m>1$}
\label{buch:ratint:bernoulli:eqn:quadrpotnenner}
\end{equation}
mit quadratischen Nennern integrieren.
Dies ist mit wohlbekannten Methoden der klassischen Analysis 
möglich, wir geben hier nur die Resultate an, die man zum Beispiel
in \cite[section 2.1]{buch:bronstein} dargestellt finden kann.

Für die Nennerpotenz $m=1$ findet man die Formel
\begin{align*}
\int 
\frac{Bz+C}{z^2+\beta z+\gamma}
\,dz
&=
\frac{B}{2} \log( z^2+\beta z+\gamma )
+
\frac{2C-\beta B}{\!\sqrt{4\gamma-\beta^2}}
\arctan\frac{2z+\beta}{\!\sqrt{4\gamma-\beta^2}}.
\end{align*}
Hier treten logarithmische Terme und ein Arkustangensterm auf.
Die Arkustangensfunktion kann ebenfalls durch Logarithmusfunktionen
in der Form $\arctan(z) =-\frac{i}2\log(z+i) +\frac{i}2\log(z-i)$
ausgedrückt werden.
Im vorliegenden Fall ist dies besonders naheliegend, weil die
beiden Nullstellen des Nenner konjugiert komplex sind.

Höhere Nennerpotenzen $m>1$ können mit der folgenden Rekursionsformel
reduziert werden:
\begin{align}
\int 
\frac{Bz+C}{(z^2+\beta z+\gamma)^m}
\,dz
&=
\frac{
(2C-\beta B)z+\beta C-2\gamma B
}{
(m-1)(4\gamma-\beta^2)(z^2+bz+c)^{m-1}
}
\notag
\\
&\qquad
+
\int \frac{
(2m-3)(2C-\beta B)
}{
(m-1)(4\gamma-\beta^2)(z^2+\beta z+\gamma)^{m-1}
}
\,dz.
\label{buch:ratint:reell:eqn:rekursion}
\end{align}
Die Rekursionsformel basiert im wesentlichen auf partieller Integration.
Sie wird sich im Abschnitt~\ref{buch:ratint:section:hermite}
als ein Spezialfall der Methode von Hermite herausstellen.

%
% Die allgemeine Form der Stammfunktion
%
\subsubsection{Die allgemeine Form der Stammfunktion}
Die vorangegangenen Berechnungen zeigen, dass das Integral einer
rationalen Funktion $f(z)$ sich immer als eine Summe von Termen
der folgenden Form darstellen lassen.
\begin{enumerate}
\item
Für jeden Linearfaktor $z-\alpha$ sind rationale Terme der Form
\[
\frac{A}{(z-\alpha)^m}
\]
möglich, wobei $m$ nicht grösser sein kann als die Potenz, mit der
der Faktor im Nenner von $f(z)$ vorkommt.

Ausserdem ist ein Term der Form
\[
A \log (z-\alpha)
\]
möglich.
\item
Für jeden quadratischen Faktor der Form $q(z)=z^2+\beta z+\gamma$
sind rationale Terme der Form $t(z) / q(z)^m$ möglich, wobei $t(z)$
ein lineares Polynom ist und $m$ nicht grösser sein kann als die
Potenz, in der $q(z)$ im Nenner von $f(z)$ vorkommt.

Ausserdem sind Terme der Form
\[
\log q(z)
=
\log (z^2 +\beta z + \gamma)
\]
und
\[
\arctan \zeta(z)
=
-\frac{i}{2}\log (\zeta+i)
+\frac{i}{2}\log (\zeta-i)
\qquad\text{mit }
\zeta(z)=\frac{2z+\beta}{\!\sqrt{4\gamma-\beta^2}}
\]
möglich.
Die Stammfunktion lässt sich also auch in diesem Fall durch
rationale Funktionen und eine Linearkombination von drei Logarithmen
ausdrücken.
Dies ist in Übereinstimmung mit den Erkenntnissen des Liouville-Prinzips,
welches in Abschnitt~\ref{buch:intalg:section:liouville} präsentiert wird.
\end{enumerate}

Die Methode von Bernoulli postuliert, dass die Argumente der
Logarithmusfunktionen die Linearfaktoren der Nullstellen des
quadratischen Nenners sind.
Es ist nicht offensichtlich, dass die Argument $\zeta+i$ und $\zeta-i$
tatsächlich solche Linearfaktoren des Polynoms $q(z)$ sind.
Wir multiplizieren die Faktoren daher aus und erhalten
\begin{align*}
(\zeta+i)(\zeta-i)
&=
\zeta^2+1
\\
&=
\frac{4z^2+4z\beta+\beta^2}{4\gamma-\beta^2}
+
1
=
\frac{4z^2+4z\beta+\beta^2+4\gamma-\beta^2}{4\gamma-\beta^2}
=
\frac{4z^2+4z\beta+4\gamma}{4\gamma-\beta^2}
\\
&=
\frac{4}{4\gamma-\beta^2}
(z^2+\beta z+\gamma).
\end{align*}
Die Faktoren $\zeta+i$ und $\zeta-i$ haben daher die gleichen
Nullstellen wie das Polynom $q(z)=z^2+\beta z+\gamma$.
Tatsächlich ist
\begin{align*}
0&=\zeta\pm i
=
\frac{2z+\beta\pm i\!\sqrt{4\gamma-\beta^2}}{\!\sqrt{4\gamma-\beta^2}}
=
\frac{2z+\beta\pm \!\sqrt{\beta^2 - 4\gamma}}{\!\sqrt{4\gamma-\beta^2}}
&&\Rightarrow&
z
&=
-\frac{\beta}2\mp \!\sqrt{\frac{\beta^2}{4}-\gamma}.
\end{align*}
Dies ist die Lösungsformel für die quadratische Gleichung
$z^2+\beta z+\gamma=0$.





